\chapter{The reactive primitives, previewed}
\label{ch:reactivity-preview}

\begin{aside}
Most of \maya{} is plain functions. A handful of places ---
animations, watched files, derived state --- want \emph{reactive
primitives}: signals that notify subscribers when their value
changes. \maya{}'s signal system lives in
\file{maya/include/maya/core/signal.hpp} and gets a full
treatment in Chapter~\ref{ch:signals}. This chapter is the
architectural preview: what they are, why they exist, where the
Elm architecture leaves room for them.
\end{aside}

\section{Signals in one paragraph}

A \code{Signal<T>} is a value that notifies subscribers when it
changes. A \code{Computed<T>} is a derived signal whose value is
a function of other signals. An \code{Effect} is a callable that
re-runs whenever one of the signals it reads changes.

If you have used SolidJS, MobX, Vue, Knockout, or Svelte 5, this
is the same machinery. The mechanics differ; the model is the
same.

\section{Why signals at all?}

Pure \code{update} works for most of an app. But two patterns
are awkward without reactive helpers:

\begin{itemize}[leftmargin=*]
  \item \textbf{Animations.} A counter that ticks once per
    frame, a loading spinner, a smooth scroll. The state is not
    business logic; it is presentation.
  \item \textbf{Derived state.} A search box that filters a
    list. The filter result is a function of the list and the
    query; recomputing it inside \code{view} every frame is
    wasteful.
\end{itemize}

You could express both with \code{Cmd} and \code{Sub}, and many
\maya{} apps do. Signals are an alternative for fine-grained
reactivity without piping every change through the message
variant.

\subsection*{The trade}

\code{update} is pure and testable. Signals introduce mutable
shared state. That is a real trade --- but it is bounded: the
signal graph is small, its subscribers are explicit, and the
mutation is observable.

For the parts of an app where the signal model is a better fit
(animations, computed columns, file watches), the cost is worth
it. For business logic, prefer \code{update}.

\begin{warning}
The two models can mix in one app, but in any given subsystem
pick one. Mixing them inside a single widget is a recipe for
``why did this not re-render?'' bugs.
\end{warning}

\section{The basic shape}

\begin{lstlisting}
auto count = Signal{0};
auto doubled = Computed{[&] { return count() * 2; }};

Effect{[&] {
    std::cout << \"count=\" << count() << \" doubled=\" << doubled() << \"\\n\";
}};

count.set(5);   // prints \"count=5 doubled=10\"
count.set(7);   // prints \"count=7 doubled=14\"
\end{lstlisting}

Three pieces:

\begin{itemize}[leftmargin=*]
  \item A \code{Signal<int>} holds a value and tracks subscribers.
  \item A \code{Computed<int>} is a function over signals. When
    its inputs change, it recomputes; subscribers of the computed
    are notified.
  \item An \code{Effect} runs once at registration, then again
    whenever any signal it read changes.
\end{itemize}

Subscriber tracking is automatic. \code{count()} inside the
effect registers a dependency; \code{count.set(\ldots)} triggers
re-run.

\section{Why this works without explicit subscribe}

The key trick: while an effect or computed is evaluating,
\maya{} sets a thread-local current node. Any \code{Signal::get}
records the current node as a subscriber. When the effect
finishes, the dependency set is fixed until the next run.

\begin{lstlisting}
template <class T>
T Signal<T>::get() {
    if (auto* node = detail::current_node) {
        subscribers_.insert(node);
        node->dependencies.insert(this);
    }
    return value_;
}
\end{lstlisting}

The dependency graph builds itself. No \code{addEventListener},
no \code{watch} declarations. The function calls themselves
register the wiring.

\subsection*{What gets tracked}

Only reads through \code{Signal::get()} (which the call operator
forwards to). Plain variable reads, function calls, and
non-signal state are invisible to the tracker.

This means: if you compute a value from a signal and a plain
variable, only the signal is tracked. Changing the plain variable
will not re-run the effect; you must read it through a signal
or wrap it in one.

\section{When to use signals in maya}

\subsection*{Animations}

A spinner is the canonical case:

\begin{lstlisting}
auto frame = Signal{0};

// somewhere: a 60Hz timer increments frame.
schedule_every(16ms, [&] { frame.set(frame() + 1); });

// in view:
Element view(const Model& m) {
    return text(spinner_frames[frame() % spinner_frames.size()]);
}
\end{lstlisting}

The view reads \code{frame()} so it auto-subscribes. The signal
changes; \maya{}'s render pipeline knows the affected sub-tree
must re-render. The model is not involved.

\subsection*{Computed columns}

A list with a filter:

\begin{lstlisting}
auto items = Signal<std::vector<Item>>{};
auto query = Signal<std::string>{};
auto filtered = Computed{[&] {
    std::vector<Item> out;
    for (auto& i : items()) {
        if (i.label.contains(query())) out.push_back(i);
    }
    return out;
}};
\end{lstlisting}

\code{filtered} recomputes only when \code{items} or
\code{query} changes. The view reads \code{filtered()};
intermediate frames see the same vector without recomputing.

\subsection*{File watches}

A signal whose value is the contents of a file, automatically
updated when the file is modified:

\begin{lstlisting}
auto config = file_signal(\"~/.config/myapp.toml\");
auto parsed = Computed{[&] { return parse_toml(config()); }};
\end{lstlisting}

Behind the scenes, \code{file\_signal} uses inotify (Linux),
kqueue (macOS), or ReadDirectoryChangesW (Windows) to watch the
file. When it changes, \code{config} updates, \code{parsed}
recomputes, every subscriber re-renders.

\section{What signals do not replace}

The Elm core stays. \code{update} is still the place for
business logic. Signals supplement it for the cases above; they
do not subsume it.

A useful rule: if the state is something a user might want to
\emph{see in a debugger}, put it in \code{Model}. If the state
is a synthetic computation (animation frame, derived index),
\code{Signal} is fine.

\section{The cost}

Signals add:

\begin{itemize}[leftmargin=*]
  \item Heap allocations: each \code{Signal} stores a
    \code{std::vector} of subscribers.
  \item Thread-local state: the dependency tracker.
  \item Mutable state: signals are by definition not pure.
  \item Glitches if you do not understand the model: re-running
    an effect can trigger another effect, which can trigger
    another\ldots
\end{itemize}

Used sparingly (a handful of signals per app), the cost is
negligible. Used as the primary state-management tool, the cost
adds up --- and you have lost the testability of pure
\code{update}.

\section*{Workshop}

\begin{enumerate}[leftmargin=*]
  \item Create a \code{Signal<int>} and an \code{Effect} that
    prints the value. Set the signal three times; observe the
    effect runs four times (once on registration, three on
    changes).
  \item Add a \code{Computed} that doubles the signal. Confirm
    the effect prints both the signal and the computed.
  \item Read \file{maya/include/maya/core/signal.hpp}. Identify
    the data members of \code{Signal<T>} and the locks (if any)
    that protect them.
\end{enumerate}

\begin{recap}
Signals, computeds, and effects form a fine-grained reactive
system. They build a dependency graph automatically through
\code{get} calls, re-run effects when inputs change, and
complement (rather than replace) the Elm architecture.
Chapter~\ref{ch:signals} treats them in full; this chapter
explains where they fit.
\end{recap}
